\documentclass{amsart}
% For dealing with references we use the comment environment
\usepackage{verbatim}
\newenvironment{reference}{\comment}{\endcomment}
%\newenvironment{reference}{}{}
\newenvironment{slogan}{\comment}{\endcomment}
\newenvironment{history}{\comment}{\endcomment}

% For commutative diagrams we use Xy-pic
\usepackage[all]{xy}

% We use 2cell for 2-commutative diagrams.
\xyoption{2cell}
\UseAllTwocells

% We use multicol for the list of chapters between chapters
\usepackage{multicol}

% This is generally recommended for better output
\usepackage{lmodern}
\usepackage[T1]{fontenc}

% For cross-file-references

% Package for hypertext links:
\usepackage{hyperref}

% For any local file, say "hello.tex" you want to link to please

% Theorem environments.
%
\theoremstyle{plain}
\newtheorem{theorem}[subsection]{Theorem}
\newtheorem{proposition}[subsection]{Proposition}
\newtheorem{lemma}[subsection]{Lemma}

\theoremstyle{definition}
\newtheorem{definition}[subsection]{Definition}
\newtheorem{example}[subsection]{Example}
\newtheorem{exercise}[subsection]{Exercise}
\newtheorem{situation}[subsection]{Situation}

\theoremstyle{remark}
\newtheorem{remark}[subsection]{Remark}
\newtheorem{remarks}[subsection]{Remarks}

\numberwithin{equation}{subsection}

% Macros
%
\def\lim{\mathop{\mathrm{lim}}\nolimits}
\def\colim{\mathop{\mathrm{colim}}\nolimits}
\def\Spec{\mathop{\mathrm{Spec}}}
\def\Hom{\mathop{\mathrm{Hom}}\nolimits}
\def\Ext{\mathop{\mathrm{Ext}}\nolimits}
\def\SheafHom{\mathop{\mathcal{H}\!\mathit{om}}\nolimits}
\def\SheafExt{\mathop{\mathcal{E}\!\mathit{xt}}\nolimits}
\def\Sch{\mathit{Sch}}
\def\Mor{\mathop{\mathrm{Mor}}\nolimits}
\def\Ob{\mathop{\mathrm{Ob}}\nolimits}
\def\Sh{\mathop{\mathit{Sh}}\nolimits}
\def\NL{\mathop{N\!L}\nolimits}
\def\CH{\mathop{\mathrm{CH}}\nolimits}
\def\proetale{{pro\text{-}\acute{e}tale}}
\def\etale{{\acute{e}tale}}
\def\QCoh{\mathit{QCoh}}
\def\Ker{\mathop{\mathrm{Ker}}}
\def\Im{\mathop{\mathrm{Im}}}
\def\Coker{\mathop{\mathrm{Coker}}}
\def\Coim{\mathop{\mathrm{Coim}}}

% Boxtimes
%
\DeclareMathSymbol{\boxtimes}{\mathbin}{AMSa}{"02}

%
% Macros for moduli stacks/spaces
%
\def\QCohstack{\mathcal{QC}\!\mathit{oh}}
\def\Cohstack{\mathcal{C}\!\mathit{oh}}
\def\Spacesstack{\mathcal{S}\!\mathit{paces}}
\def\Quotfunctor{\mathrm{Quot}}
\def\Hilbfunctor{\mathrm{Hilb}}
\def\Curvesstack{\mathcal{C}\!\mathit{urves}}
\def\Polarizedstack{\mathcal{P}\!\mathit{olarized}}
\def\Complexesstack{\mathcal{C}\!\mathit{omplexes}}
% \Pic is the operator that assigns to X its picard group, usage \Pic(X)
% \Picardstack_{X/B} denotes the Picard stack of X over B
% \Picardfunctor_{X/B} denotes the Picard functor of X over B
\def\Pic{\mathop{\mathrm{Pic}}\nolimits}
\def\Picardstack{\mathcal{P}\!\mathit{ic}}
\def\Picardfunctor{\mathrm{Pic}}
\def\Deformationcategory{\mathcal{D}\!\mathit{ef}}

\setlength{\emergencystretch}{3em}
\begin{document}
\title[Stacks Project, Tag 0E2F]{252 --- Existence and arbitrary base change of relative dualizing complexes for flat finitely presented algebras}
\maketitle
\noindent Copyright (C) 2005--2025 Johan de Jong. Stacks Project authors. Source: \href{https://stacks.math.columbia.edu/tag/0E2F}{Tag 0E2F}.

\noindent Editorial difficulty note (not author text). Difficulty H2: The proof first descends the algebra to a finite-type integer model. For arbitrary non-Noetherian base changes it builds a relatively flat Hom model, proves universal injectivity and flatness of its successive cokernels by residue-field testing, and identifies degree-zero cohomology compatibly with the rigidifying generator. These are the core construction and base-change obligations..

\noindent Editorial provenance note (not author text). History: excerpt from revision \texttt{a0444\allowbreak 6e57e\allowbreak c1fbc\allowbreak 252a8\allowbreak 71afc\allowbreak ec775\allowbreak 2fb28\allowbreak 07b14}, dualizing.tex, lines 5473--5578. Reference presentation was adapted; original-author context and core support blocks were retained or added. The mathematical wording is unchanged. Licensed under GNU FDL 1.2 or later; no invariant sections or cover texts. See ../licenses/COPYING. Context blocks reproduce author definitions and supporting results; they are not additional counted problems.

\begin{definition}
\label{definition-relative-dualizing-complex}
Let $R \to A$ be a flat ring map of finite presentation.
A {\it relative dualizing complex} is an object $K \in D(A)$ such that
\begin{enumerate}
\item $K$ is $R$-perfect (More on Algebra, Definition
\ref{more-algebra-definition-relatively-perfect}), and
\item $R\Hom_{A \otimes_R A}(A, K \otimes_A^\mathbf{L} (A \otimes_R A))$
is isomorphic to $A$.
\end{enumerate}
\end{definition}

\begin{definition}
\label{more-algebra-definition-relatively-perfect}
Let $R \to A$ be a flat ring map of finite presentation.
An object $K$ of $D(A)$ is {\it $R$-perfect} or {\it perfect relative to $R$}
if $K$ is pseudo-coherent
(Definition \ref{more-algebra-definition-pseudo-coherent})
and has finite tor dimension over $R$
(Definition \ref{more-algebra-definition-tor-amplitude}).
\end{definition}

\begin{definition}
\label{more-algebra-definition-pseudo-coherent}
Let $R$ be a ring. Denote $D(R)$ its derived category.
Let $m \in \mathbf{Z}$.
\begin{enumerate}
\item An object $K^\bullet$ of $D(R)$ is {\it $m$-pseudo-coherent}
if there exists a bounded complex $E^\bullet$ of finite free $R$-modules
and a morphism $\alpha : E^\bullet \to K^\bullet$ such that
$H^i(\alpha)$ is an isomorphism for $i > m$ and $H^m(\alpha)$
is surjective.
\item An object $K^\bullet$ of $D(R)$ is {\it pseudo-coherent}
if it is quasi-isomorphic to a bounded above complex of finite
free $R$-modules.
\item An $R$-module $M$ is called {\it $m$-pseudo-coherent}
if $M[0]$ is an $m$-pseudo-coherent object of $D(R)$.
\item An $R$-module $M$ is called
{\it pseudo-coherent}\footnote{This clashes with what is meant by
a pseudo-coherent module in \href{https://stacks.math.columbia.edu/bibliography}{Bibliography: Bourbaki-CA}.}
if $M[0]$ is a pseudo-coherent object of $D(R)$.
\end{enumerate}
\end{definition}

\begin{definition}
\label{more-algebra-definition-tor-amplitude}
Let $R$ be a ring. Denote $D(R)$ its derived category.
Let $a, b \in \mathbf{Z}$.
\begin{enumerate}
\item An object $K^\bullet$ of $D(R)$ has
{\it tor-amplitude in $[a, b]$}
if $H^i(K^\bullet \otimes_R^\mathbf{L} M) = 0$ for all $R$-modules
$M$ and all $i \not \in [a, b]$.
\item An object $K^\bullet$ of $D(R)$ has {\it finite tor dimension}
if it has tor-amplitude in $[a, b]$ for some $a, b$.
\item An $R$-module $M$ has {\it tor dimension $\leq d$}
if $M[0]$ as an object of $D(R)$ has tor-amplitude in $[-d, 0]$.
\item An $R$-module $M$ has {\it finite tor dimension}
if $M[0]$ as an object of $D(R)$ has finite tor dimension.
\end{enumerate}
\end{definition}

\begin{lemma}
\label{lemma-base-change-relative-dualizing}
Let $R \to A$ be a flat ring map of finite presentation. Then
\begin{enumerate}
\item there exists a relative dualizing complex $K$ in $D(A)$, and
\item for any ring map $R \to R'$ setting $A' = A \otimes_R R'$
and $K' = K \otimes_A^\mathbf{L} A'$, then $K'$ is a
relative dualizing complex for $R' \to A'$.
\end{enumerate}
Moreover, if
$$
\xi : A \longrightarrow K \otimes_A^\mathbf{L} (A \otimes_R A)
$$
is a generator for the cyclic module
$\Hom_{D(A \otimes_R A)}(A, K \otimes_A^\mathbf{L} (A \otimes_R A))$
then in (2) the derived base change of $\xi$ by
$A \otimes_R A \to A' \otimes_{R'} A'$ is a generator for
the cyclic module
$\Hom_{D(A' \otimes_{R'} A')}(A',
K' \otimes_{A'}^\mathbf{L} (A' \otimes_{R'} A'))$
\end{lemma}

\begin{proof}
We first reduce to the Noetherian case. By
Algebra, Lemma \href{https://stacks.math.columbia.edu/tag/02JO}{lemma flat finite presentation limit flat (Tag 02JO)}
there exists a finite type $\mathbf{Z}$ subalgebra $R_0 \subset R$
and a flat finite type ring map $R_0 \to A_0$ such that
$A = A_0 \otimes_{R_0} R$. By Lemma \href{https://stacks.math.columbia.edu/tag/0E2E}{lemma relative dualizing noetherian (Tag 0E2E)}
there exists a relative
dualizing complex $K_0 \in D(A_0)$.
Thus if we show (2) for $K_0$, then we find that
$K_0 \otimes_{A_0}^\mathbf{L} A$ is
a dualizing complex for $R \to A$ and that it also satisfies (2)
by transitivity of derived base change.
The uniqueness of relative dualizing complexes
(Lemma \href{https://stacks.math.columbia.edu/tag/0E2D}{lemma uniqueness relative dualizing (Tag 0E2D)})
then shows that this holds for
any relative dualizing complex.

\medskip\noindent
Assume $R$ Noetherian and let $K$ be a relative dualizing complex
for $R \to A$. Given a ring map $R \to R'$ set $A' = A \otimes_R R'$
and $K' = K \otimes_A^\mathbf{L} A'$. To finish the proof we have
to show that $K'$ is a relative dualizing complex for $R' \to A'$.
By More on Algebra, Lemma
\href{https://stacks.math.columbia.edu/tag/0DHW}{lemma base change relatively perfect (Tag 0DHW)}
we see that $K'$ is $R'$-perfect in all cases.
By Lemmas \href{https://stacks.math.columbia.edu/tag/0BZV}{lemma base change relative algebraic (Tag 0BZV)} and
\href{https://stacks.math.columbia.edu/tag/0E2E}{lemma relative dualizing noetherian (Tag 0E2E)}
if $R'$ is Noetherian, then $K'$ is a relative dualizing complex
for $R' \to A'$ (in either sense).
Transitivity of derived tensor product shows that
$K \otimes_A^\mathbf{L} (A \otimes_R A)
\otimes_{A \otimes_R A}^\mathbf{L} (A' \otimes_{R'} A') =
K' \otimes_{A'}^\mathbf{L} (A' \otimes_{R'} A')$.
Flatness of $R \to A$ guarantees that
$A \otimes_{A \otimes_R A}^\mathbf{L} (A' \otimes_{R'} A') = A'$;
namely $A \otimes_R A$ and $R'$ are tor independent over $R$
so we can apply More on Algebra, Lemma
\href{https://stacks.math.columbia.edu/tag/0661}{lemma base change comparison (Tag 0661)}.
Finally, $A$ is pseudo-coherent as an $A \otimes_R A$-module
by More on Algebra, Lemma
\href{https://stacks.math.columbia.edu/tag/0E1T}{lemma more relative pseudo coherent is moot (Tag 0E1T)}. Thus
we have checked all the assumptions of
More on Algebra, Lemma
\href{https://stacks.math.columbia.edu/tag/0E1U}{lemma compute RHom relatively perfect (Tag 0E1U)}.
We find there exists a bounded below complex
$E^\bullet$ of $R$-flat finitely presented $A \otimes_R A$-modules
such that $E^\bullet \otimes_R R'$ represents
$R\Hom_{A' \otimes_{R'} A'}(A',
K' \otimes_{A'}^\mathbf{L} (A' \otimes_{R'} A'))$
and these identifications are compatible with derived base change.
Let $n \in \mathbf{Z}$, $n \not = 0$.
Define $Q^n$ by the sequence
$$
E^{n - 1} \to E^n \to Q^n \to 0
$$
Since $\kappa(\mathfrak p)$ is a Noetherian ring, we know that
$H^n(E^\bullet \otimes_R \kappa(\mathfrak p)) = 0$, see remarks above.
Chasing diagrams this means that
$$
Q^n \otimes_R \kappa(\mathfrak p) \to E^{n + 1} \otimes_R \kappa(\mathfrak p)
$$
is injective. Hence for a prime $\mathfrak q$ of $A \otimes_R A$
lying over $\mathfrak p$ we have $Q^n_\mathfrak q$ is $R_\mathfrak p$-flat
and $Q^n_\mathfrak p \to E^{n + 1}_\mathfrak q$ is
$R_\mathfrak p$-universally injective, see
Algebra, Lemma \href{https://stacks.math.columbia.edu/tag/00ME}{lemma mod injective (Tag 00ME)}.
Since this holds for all primes,
we conclude that $Q^n$ is $R$-flat
and $Q^n \to E^{n + 1}$ is $R$-universally injective. In particular
$H^n(E^\bullet \otimes_R R') = 0$ for any ring map $R \to R'$.
Let $Z^0 = \Ker(E^0 \to E^1)$. Since there is an exact sequence
$0 \to Z^0 \to E^0 \to E^1 \to Q^1 \to 0$ we see that $Z^0$
is $R$-flat and that
$Z^0 \otimes_R R' = \Ker(E^0 \otimes_R R' \to E^1 \otimes_R R')$
for all $R \to R'$. Then the short exact sequence
$0 \to Q^{-1} \to Z^0 \to H^0(E^\bullet) \to 0$
shows that
$$
H^0(E^\bullet \otimes_R R') = H^0(E^\bullet) \otimes_R R'
= A \otimes_R R' = A'
$$
as desired. This equality furthermore gives the final assertion
of the lemma.
\end{proof}
\end{document}
